\section*{Correct Table 1 to the citywide descriptive sample}
\textit{September 11, 2026. Agent-drafted correction following Jacob's clarification.}

The preceding common-sample table used the wrong population. Jacob intended
all citywide new construction, not the 500-foot estimation sample. The older
\texttt{summary\_stats\_new\_construction} script confirms that its summary had
no boundary-distance cutoff. The current analysis-data producer also imposes a
1,500-foot restriction, so removing the 500-foot filter alone would not suffice.

The table now reads the full cleaned project ledger and construction-year
boundary distances directly. It retains 2006--2022 construction with both FAR
and DUPAC usable and positive units, without distance, score, zoning-control or
regression-complete-case restrictions. The two-column layout is retained.
Building-type classification follows the current reviewed classification and is
checked against the existing analysis records where they overlap. Regression
samples and methods are unchanged.

The citywide common descriptive sample has 13,635 observations, including 2,582
multifamily observations. Median distance is 914.94 feet overall and 825.93 feet
for multifamily. Mean FAR is 1.158 and 2.085; mean DUPAC is 26.936 and 65.181;
mean units are 5.805 and 26.105. The first citywide version omitted the
ward-pair/segment-count row. The note and accompanying prose explicitly state
the citywide scope. This entry supersedes the prior entry's table population.

Jacob subsequently chose the September 3 draft as the presentation reference,
with citywide Table 1 as the exception. Its original caption, two columns, row
order and dividing rules are restored, including the ward/segment row. That
row reports 139 nearest ward pairs and 1,388 segments overall, and 112 pairs
and 805 segments for multifamily. The note clarifies that these are ward pairs,
not a count of Chicago wards. Every eligible citywide project is assigned to
the nearest segment of its recorded construction-year ward boundary, without
a distance cutoff; assignments agree with the existing analysis records where
both are available. Other descriptive statistics are unchanged.

The manuscript comparison uses revision \texttt{2e406330}. Outside the necessary
citywide sample description, differences in wording are limited to the numerical
updates and Jacob's previously approved permit-period and placebo sentences.
The original significance and donut descriptions remain. Figure layouts and
labels are unchanged; plotted values are the rerun results. The July draft's
extra stories and stringency rows and older building-type definition are not
adopted. This restores presentation without undoing the approved data corrections
or outcome-specific regression samples.

\smallskip
\noindent\textit{Reproduction:} Run Make in
\texttt{tasks/density\_main\_results/code/}, then in \texttt{paper/}.
The summary CSV now produces the table and data report through separate file
rules. The saved report has two rows and eight columns, adding the two geographic
counts. Builds and rendered exhibits were checked in the completed replication
clone on the \texttt{conference-replication} branch.

\smallskip
\noindent\textit{Approved wording revision.} Jacob approved the displayed
before-and-after wording changes across the paper. Table 1 is now titled
``New Residential Construction in Chicago, 2006--2022.'' Its note defines the
measurements and multifamily buildings directly. Row labels identify units per
project, ward pairs and boundary segments, and the number of projects. The
citywide common sample and existing building-type classification are retained.
The approved edits also simplify the introduction, data descriptions, and
appendix table titles and notes, and remove unnecessary ``final'' references
to stringency scores. Jacob reserved the three factual passages on density
eligibility, the event-study period, and the count of manually reviewed projects
for a separate discussion. Those passages are unchanged in this revision.
